% ============================================================
%  Übungsaufgaben Template (my_dhbw Stil, uebung_*.tex)
%  Header "Übungsblatt", grün eingefärbte <details>-Lösungen.
%  Renderer: pdflatex (zweimal)
% ============================================================
\documentclass[11pt, a4paper]{article}

\usepackage[ngerman]{babel}
\usepackage{fontspec}
\setmainfont[
  Path=/usr/share/fonts/TTF/,
  Extension=.ttf,
  BoldFont=DejaVuSans-Bold,
  ItalicFont=DejaVuSans-Oblique,
  BoldItalicFont=DejaVuSans-BoldOblique
]{DejaVuSans}
\setsansfont[
  Path=/usr/share/fonts/TTF/,
  Extension=.ttf,
  BoldFont=DejaVuSans-Bold,
  ItalicFont=DejaVuSans-Oblique,
  BoldItalicFont=DejaVuSans-BoldOblique
]{DejaVuSans}
\setmonofont[Path=/usr/share/fonts/TTF/,Extension=.ttf]{DejaVuSansMono}
\usepackage[top=2cm, bottom=2cm, left=2.4cm, right=2.4cm, footskip=1cm]{geometry}
\usepackage{amsmath, amssymb}
\usepackage{xcolor}
\usepackage{enumitem}
\usepackage{fancyhdr}
\usepackage{lastpage}
\usepackage{tabularx}
\usepackage{booktabs}
\usepackage{microtype}
\usepackage{parskip}

\usepackage{array}
\usepackage{longtable}
\usepackage{ltablex}
\keepXColumns
\usepackage{colortbl}
\usepackage[most,breakable]{tcolorbox}
\usepackage{listings}
\usepackage[normalem]{ulem}
\usepackage{pdflscape}
\usepackage{titlesec}
\usepackage[colorlinks=true, linkcolor=accent, urlcolor=accent]{hyperref}

\renewcommand{\familydefault}{\sfdefault}

% ── Theme ─────────────────────────────────────────────────────
\definecolor{accent}{RGB}{0, 60, 113}
\definecolor{primaryblue}{RGB}{0, 60, 113}
\definecolor{loesung}{RGB}{0, 100, 0}
\definecolor{codebg}{RGB}{245, 245, 245}
\definecolor{figurebg}{RGB}{232, 241, 255}
\definecolor{tablehintbg}{RGB}{240, 255, 240}
\definecolor{solutionbg}{RGB}{240, 252, 240}
\definecolor{rulegray}{RGB}{180, 180, 180}
\definecolor{tableheadbg}{RGB}{0, 60, 113}

% ── Header/Footer ─────────────────────────────────────────────
\pagestyle{fancy}
\fancyhf{}
\renewcommand{\headrulewidth}{0.4pt}
\fancyhead[L]{\footnotesize\color{accent}\nichttitle}
\fancyhead[R]{\footnotesize\color{accent}\nichtsubtitle}
\fancyfoot[R]{\footnotesize\color{accent}Blatt \thepage\,/\,\pageref{LastPage}}

% ── Typografie ────────────────────────────────────────────────
\titleformat{\section}
    {\color{accent}\large\bfseries}{}{0pt}{}
    [\vspace{-4pt}\color{accent}\rule{\linewidth}{1.5pt}]
\titleformat{\subsection}
    {\normalsize\bfseries\color{accent!85!black}}{}{0pt}{}{}
\titleformat{\subsubsection}
    {\small\bfseries\color{accent!70!black}}{}{0pt}{}{}
\titlespacing*{\section}{0pt}{10pt}{3pt}
\titlespacing*{\subsection}{0pt}{6pt}{2pt}
\titlespacing*{\subsubsection}{0pt}{4pt}{1pt}

\setlength{\parindent}{0pt}
\setlength{\parskip}{6pt}
\renewcommand{\arraystretch}{1.2}
\setlist[itemize]{noitemsep, topsep=3pt, parsep=0pt, leftmargin=1.4em}
\setlist[enumerate]{noitemsep, topsep=3pt, parsep=0pt, leftmargin=1.6em}

% ── Boxen: solutionbox = Lösungsbox in grün ──────────────────
\newtcolorbox{figurebox}[1]{
  colback=figurebg, colframe=accent!50,
  title={Abbildung: #1}, fonttitle=\small\bfseries,
  breakable, before skip=6pt, after skip=6pt
}
\newtcolorbox{tablehintbox}[1]{
  colback=tablehintbg, colframe=green!50!black!50,
  title={Tabelle: #1}, fonttitle=\small\bfseries,
  breakable, before skip=6pt, after skip=6pt
}
\newtcolorbox{solutionbox}[1]{
  enhanced,
  colback=solutionbg, colframe=loesung,
  title={#1}, fonttitle=\small\bfseries,
  coltitle=white, colbacktitle=loesung,
  breakable, before skip=8pt, after skip=8pt
}

\lstset{
  basicstyle=\ttfamily\small,
  backgroundcolor=\color{codebg},
  breaklines=true,
  frame=single, framerule=0pt,
  xleftmargin=4pt, xrightmargin=4pt,
  aboveskip=6pt, belowskip=6pt,
  keepspaces=true
}

\providecommand{\nichttitle}{Übungsblatt}
\providecommand{\nichtsubtitle}{}

\renewcommand{\nichttitle}{Analysis 2}
\renewcommand{\nichtsubtitle}{Übungsblatt 3}


\begin{document}

\begin{center}
{\Large\bfseries\color{accent}\nichttitle}
\ifx\nichtsubtitle\empty\else
  \\[2pt]{\normalsize\color{accent!80!black}\nichtsubtitle}
\fi
\\[2pt]
\color{accent}\rule{0.4\linewidth}{0.8pt}\color{black}
\end{center}
\vspace{4pt}

% ── BODY ──────────────────────────────────────────────────────

\section{Übungsblatt 3 — Vektoranalysis}


\noindent\textcolor{rulegray}{\rule{\linewidth}{0.4pt}}


\paragraph{Aufgabe 1 — Totales Differenzial \& Richtungsableitung \textit{(15 Punkte)}}\mbox{}\\[2pt]

Betrachte die Funktion $f(x,y) = x^2 \cdot \ln(y)$.


a) (3 P) Klassifiziere $f$ — handelt es sich um ein Skalar- oder Vektorfeld? Begründe anhand der Definition.


b) (4 P) Bestimme das totale Differenzial $df$ und werte es an der Stelle $(1, e)$ in den Achsenrichtungen $(dx=1, dy=0)$ und $(dx=0, dy=1)$ aus.


c) (8 P) Berechne die Richtungsableitung an $(1, e)$ in der Richtung, die mit der $x$-Achse den Winkel \$60°$ einschließt (also $dy = \textbackslash{}sqrt\{3\}\textbackslash{}cdot dx$). Bestimme dazu zuerst den korrekten normierten Richtungsvektor $\textbackslash{}vec\{a\}$ und begründe, warum die Normierung $|\textbackslash{}vec\{a\}|=1\$ zwingend nötig ist.


\textbf{Lösung:}


a) $f: \mathbb{R}^2 \to \mathbb{R}$ — jedem Punkt $(x,y)$ wird genau eine reelle Zahl zugeordnet. Also ein \textbf{Skalarfeld}.


b) Partielle Ableitungen: $\frac{\partial f}{\partial x} = 2x\ln(y)$, $\frac{\partial f}{\partial y} = \frac{x^2}{y}$.

\[
df = 2x\ln(y)\,dx + \frac{x^2}{y}\,dy
\]

An $(1,e)$: $\frac{\partial f}{\partial x} = 2\cdot 1\cdot \ln(e) = 2$ und $\frac{\partial f}{\partial y} = \frac{1}{e}$.

Achsenrichtung $x$: $df = 2$. Achsenrichtung $y$: $df = 1/e$.


c) Allgemeiner Richtungsvektor: $\vec{v} = (1, \sqrt{3})$, Betrag $|\vec{v}| = \sqrt{1+3} = 2$. Normiert:

\[
\vec{a} = \tfrac{1}{2}(1, \sqrt{3}) = (0{,}5;\ 0{,}5\sqrt{3})
\]

Probe: $|\vec{a}|^2 = 0{,}25 + 0{,}75 = 1$ ✓.

\[
\frac{\partial f}{\partial \vec{a}}\bigg|_{(1,e)} = 2\cdot 0{,}5 + \frac{1}{e}\cdot 0{,}5\sqrt{3} = 1 + \frac{\sqrt{3}}{2e}
\]

Begründung Normierung: Die Richtungsableitung ist als Skalarprodukt $\nabla f \cdot \vec{a}$ definiert. Wäre $|\vec{a}| \neq 1$, würde das Ergebnis mit der Länge des Vektors skalieren — ein "doppelt so langer" Vektor in dieselbe Richtung lieferte das doppelte Ergebnis. Sie soll aber nur die Änderungsrate \textbf{pro Längeneinheit} in einer Richtung beschreiben, deshalb fordert die Definition $|\vec{a}|=1$.



\noindent\textcolor{rulegray}{\rule{\linewidth}{0.4pt}}


\paragraph{Aufgabe 2 — Gradient, kritische Punkte und steilster Anstieg \textit{(15 Punkte)}}\mbox{}\\[2pt]

Gegeben: $f(x,y) = 2 - x - 2y + xy$.


a) (4 P) Bestimme $\operatorname{grad}(f)$.


b) (4 P) Finde alle kritischen Punkte (notwendige Bedingung für Extremstellen).


c) (7 P) Ein Wanderer steht im Punkt $P=(3,2)$ auf der durch $f$ beschriebenen Höhenkarte. In welche Richtung muss er gehen, damit er den steilsten Aufstieg nimmt? Wie groß ist die momentane Anstiegsrate? Gib die Aufstiegsrichtung als normierten Vektor an.


\textbf{Lösung:}


a) $\nabla f = \begin{pmatrix} -1+y \\ -2+x \end{pmatrix}$.


b) $\nabla f = \vec{0}$ liefert das Gleichungssystem $-1+y=0$ und $-2+x=0$, also $\boxed{(x,y) = (2,1)}$ — ein einziger kritischer Punkt.


c) Bei $P=(3,2)$:

\[
\nabla f(3,2) = \begin{pmatrix} -1+2 \\ -2+3 \end{pmatrix} = \begin{pmatrix} 1 \\ 1 \end{pmatrix}
\]

Eigenschaft des Gradienten: er zeigt in Richtung des steilsten Anstiegs, sein Betrag ist die Anstiegsrate.


Anstiegsrate: $|\nabla f| = \sqrt{1^2 + 1^2} = \sqrt{2}$.


Normierte Aufstiegsrichtung: $\vec{e} = \frac{1}{\sqrt{2}}\begin{pmatrix}1\\1\end{pmatrix} = \begin{pmatrix}1/\sqrt{2}\\1/\sqrt{2}\end{pmatrix}$.


Der Wanderer bewegt sich also genau diagonal nach Nordosten und gewinnt dabei pro Schritt $\sqrt{2} \approx 1{,}414$ Höheneinheiten.



\noindent\textcolor{rulegray}{\rule{\linewidth}{0.4pt}}


\paragraph{Aufgabe 3 — Divergenz vs. Rotation: Analyse eines Vektorfelds \textit{(20 Punkte)}}\mbox{}\\[2pt]

Gegeben das Vektorfeld $\vec{f}(x,y) = (x^2y,\ x - y^2)^T$.


a) (4 P) Berechne $\operatorname{div}(\vec{f})$ und $\operatorname{rot}(\vec{f})$ (in $\mathbb{R}^2$ skalar).


b) (3 P) Bestimme alle Punkte/Kurven, an denen $\vec{f}$ \textbf{quellenfrei} ist.


c) (3 P) Bestimme alle Punkte/Kurven, an denen $\vec{f}$ \textbf{wirbelfrei} ist.


d) (4 P) Existieren Punkte, an denen $\vec{f}$ \textbf{gleichzeitig} quellen- und wirbelfrei ist? Falls ja, gib sie an.


e) (6 P) Vergleiche Divergenz und Rotation als Charakterisierungen eines Vektorfelds. Begründe anhand des Beispielfelds $(-y, x)^T$ aus dem Kapitel, warum aus $\operatorname{div}(\vec{f})=0$ nicht $\operatorname{rot}(\vec{f})=0$ folgt — und umgekehrt.


\textbf{Lösung:}


a) $\operatorname{div}(\vec{f}) = \frac{\partial}{\partial x}(x^2y) + \frac{\partial}{\partial y}(x-y^2) = 2xy - 2y = 2y(x-1)$.

$\operatorname{rot}(\vec{f}) = \frac{\partial}{\partial x}(x-y^2) - \frac{\partial}{\partial y}(x^2y) = 1 - x^2$.


b) Quellenfrei $\Leftrightarrow 2y(x-1) = 0 \Leftrightarrow y=0$ \textbf{oder} $x=1$ (Vereinigung zweier Geraden).


c) Wirbelfrei $\Leftrightarrow 1-x^2 = 0 \Leftrightarrow x = \pm 1$ (zwei vertikale Geraden).


d) Schnittmenge: $x=\pm 1$ und ($y=0$ oder $x=1$).

\begin{itemize}
  \item $x=1$: jede Stelle $(1,y)$ erfüllt beides → ganze Gerade $x=1$.
  \item $x=-1$: zusätzlich $y=0$ nötig → Punkt $(-1, 0)$.
\end{itemize}

Lösungsmenge: $\{(1,y)\,|\,y\in\mathbb{R}\}\ \cup\ \{(-1,0)\}$.


e) Divergenz misst die \textbf{Quellenstärke} (skalar): wie viel Materie/Fluss in einem Punkt entsteht oder verschwindet. Rotation misst die \textbf{Wirbelstärke} (skalar in $\mathbb{R}^2$): wie stark das Feld um einen Punkt rotiert. Beide Operatoren extrahieren \textbf{unabhängige} Eigenschaften.


Beispiel: $\vec{g}(x,y) = (-y, x)^T$. Es gilt $\operatorname{div}(\vec{g}) = 0 + 0 = 0$ (überall quellenfrei), aber $\operatorname{rot}(\vec{g}) = 1 - (-1) = 2 \neq 0$ (überall wirbelhaft). Umgekehrt ist das radiale Feld $(x,y)^T$ wirbelfrei ($\operatorname{rot}=0$), aber nicht quellenfrei ($\operatorname{div}=2$). Damit kann man von einer der beiden Größen nicht auf die andere schließen.



\noindent\textcolor{rulegray}{\rule{\linewidth}{0.4pt}}


\paragraph{Aufgabe 4 — Gradient Descent für lineare Regression \textit{(25 Punkte)}}\mbox{}\\[2pt]

Wir trainieren eine einfache lineare Regression $\hat{y}_i = \beta_0 + \beta_1 x_i$ auf den drei Datenpunkten:



{\small
\begin{tabularx}{\linewidth}{XXX}
\hline
\rowcolor{tableheadbg}
{\color{white}\textbf{$i$}} & {\color{white}\textbf{$x_i$}} & {\color{white}\textbf{$y_i$}} \\[2pt]
\hline
\endhead
\hline
\endfoot
1 & 1 & 2 \\
2 & 2 & 2 \\
3 & 3 & 4 \\
\end{tabularx}
}


Als Loss verwenden wir den \textbf{mittleren quadratischen Fehler} (MSE, ohne Wurzel):

\[
f(\beta_0, \beta_1) = \frac{1}{3}\sum_{i=1}^{3}(y_i - \beta_0 - \beta_1 x_i)^2
\]


a) (5 P) Bestimme $\nabla f(\beta_0, \beta_1)$ in geschlossener Form.


b) (8 P) Führe ausgehend vom Startwert $\beta^{(0)} = (0, 0)$ \textbf{einen} Gradient-Descent-Schritt mit Lernrate $\eta = 0{,}1$ aus. Gib $\beta^{(1)}$ explizit an.


c) (6 P) Im Kapitel steht: "$\operatorname{rot}(\operatorname{grad}(f)) = 0$ → keine Endlosschleife". Erläutere unter Verwendung des Satzes von Schwarz, was diese Aussage konkret für Gradient Descent garantiert.


d) (6 P) Erkläre anhand der Stichworte "Diskretisierung" und "lokale vs. globale Minima", unter welchen Bedingungen Gradient Descent trotz korrekter Mathematik \textbf{kein} zufriedenstellendes Ergebnis liefert. Nenne zu jedem Stichwort ein konkretes Failure-Szenario.


\textbf{Lösung:}


a) Mit der Kettenregel:

\[
\frac{\partial f}{\partial \beta_0} = \frac{1}{3}\sum_{i=1}^{3} 2(y_i - \beta_0 - \beta_1 x_i)(-1) = -\frac{2}{3}\sum_{i}(y_i - \beta_0 - \beta_1 x_i)
\]

\[
\frac{\partial f}{\partial \beta_1} = -\frac{2}{3}\sum_{i} x_i(y_i - \beta_0 - \beta_1 x_i)
\]


b) Bei $\beta^{(0)}=(0,0)$ ist der Residualterm einfach $y_i$:


\begin{itemize}
  \item $\sum y_i = 2+2+4 = 8 \Rightarrow \partial_{\beta_0} f = -\frac{2}{3}\cdot 8 = -\frac{16}{3}$
  \item $\sum x_i y_i = 1\cdot 2 + 2\cdot 2 + 3\cdot 4 = 18 \Rightarrow \partial_{\beta_1} f = -\frac{2}{3}\cdot 18 = -12$
\end{itemize}

Update-Regel: $\beta^{(1)} = \beta^{(0)} - \eta \cdot \nabla f$.

\[
\beta_0^{(1)} = 0 - 0{,}1\cdot\left(-\tfrac{16}{3}\right) = \tfrac{16}{30} = \tfrac{8}{15} \approx 0{,}533
\]

\[
\beta_1^{(1)} = 0 - 0{,}1\cdot(-12) = 1{,}2
\]

\[
\boxed{\beta^{(1)} \approx (0{,}533;\ 1{,}2)}
\]


c) Der Satz von Schwarz besagt: für $f \in C^2$ gilt $\partial^2 f / \partial \beta_0 \partial \beta_1 = \partial^2 f / \partial \beta_1 \partial \beta_0$. Daraus folgt unmittelbar

\[
\operatorname{rot}(\nabla f) = \frac{\partial^2 f}{\partial \beta_0 \partial \beta_1} - \frac{\partial^2 f}{\partial \beta_1 \partial \beta_0} = 0
\]

Konsequenz: das Gradientenfeld ist \textbf{wirbelfrei} und besitzt damit ein Potential (nämlich $f$ selbst). Würde GD entlang einer geschlossenen Bahn rotieren ohne $f$ zu reduzieren, müsste die Bahn netto Rotation einschließen — dies widerspräche $\operatorname{rot}(\nabla f)=0$. Jeder Schritt entlang $-\nabla f$ verringert (für hinreichend kleines $\eta$) den Loss strikt, das Verfahren kann also nicht in einer kreisförmigen Endlosschleife mit konstantem Loss hängen bleiben.


d) \textbf{Diskretisierung}: GD nähert den kontinuierlichen Fluss $-\nabla f$ durch endliche Schritte $\eta \nabla f$. Ist $\eta$ zu groß, "überspringt" das Verfahren das Minimum und divergiert (Loss explodiert oszillierend). Ist $\eta$ zu klein, bleibt es in flachen Plateaus quasi stehen.


\textbf{Lokale vs. globale Minima}: Bei nicht-konvexem Loss (im Kapitel-Beispiel zwar nicht der Fall — MSE ist konvex — aber bei tieferen Modellen typisch) garantiert die notwendige Bedingung $\nabla f = 0$ nur einen kritischen Punkt. GD konvergiert also möglicherweise zu einem lokalen Minimum, aus dem es das globale nicht mehr unterscheiden kann.



\noindent\textcolor{rulegray}{\rule{\linewidth}{0.4pt}}


\paragraph{Aufgabe 5 — Code-Analyse: Fehler in der Richtungsableitung \textit{(15 Punkte)}}\mbox{}\\[2pt]

Ein Praktikant implementiert die Richtungsableitung von $f(x,y) = \sqrt{xy}$ an der Stelle $(4,4)$:


\begin{lstlisting}[language=python]
def richtungsableitung(f, point, v):
    # partielle Ableitungen an point
    fx = partial(f, 'x', point)   # bei (4,4): 0.5
    fy = partial(f, 'y', point)   # bei (4,4): 0.5
    return fx*v[0] + fy*v[1]

# Aufruf: richtungsableitung(f, (4,4), (1, sqrt(3)))
# Ausgabe: 0.5 + 0.5 * sqrt(3)  ≈ 1.366
\end{lstlisting}


a) (4 P) Welche Definitionsanforderung wird hier verletzt? Welcher Wert sollte laut der Beispieltabelle des Kapitels für $f=\sqrt{xy}$ an $(4,4)$ in Richtung $\vec{a}=(0{,}5;\ 0{,}5\sqrt{3})$ tatsächlich herauskommen?


b) (5 P) Korrigiere den Code (Pseudocode genügt). Erkläre die Korrektur in einem Satz.


c) (6 P) Was passiert, wenn man im \textbf{fehlerhaften} Code stattdessen $\vec{v}' = (2, 2\sqrt{3})$ übergibt — und was passiert im \textbf{korrigierten} Code? Berechne beide Ergebnisse explizit und schließe daraus, welche fundamentale Eigenschaft der Richtungsableitung der Original-Code verletzt.


\textbf{Lösung:}


a) Die Richtungsableitung ist nur für \textbf{normierte} Richtungsvektoren $|\vec{a}|=1$ definiert. Übergeben wird $\vec{v}=(1, \sqrt{3})$ mit $|\vec{v}|=\sqrt{1+3}=2 \neq 1$. Korrektes Ergebnis laut Tabelle: $0{,}25\cdot(1+\sqrt{3}) \approx 0{,}683$.


b) Korrektur: vor dem Skalarprodukt normieren.

\begin{lstlisting}[language=python]
def richtungsableitung(f, point, v):
    norm = sqrt(v[0]**2 + v[1]**2)
    a = (v[0]/norm, v[1]/norm)
    fx = partial(f, 'x', point)
    fy = partial(f, 'y', point)
    return fx*a[0] + fy*a[1]
\end{lstlisting}

Erklärung: Die Skalierung mit $1/|\vec{v}|$ stellt sicher, dass das Resultat unabhängig von der Länge des Eingabevektors nur die Änderungsrate \textbf{pro Längeneinheit in der Richtung} misst.


c) Mit $\vec{v}' = (2, 2\sqrt{3})$, $|\vec{v}'| = \sqrt{4+12} = 4$:


\textbf{Fehlerhafter Code:} $0{,}5\cdot 2 + 0{,}5\cdot 2\sqrt{3} = 1 + \sqrt{3} \approx 2{,}732$ — \textbf{doppelt} so groß wie zuvor.


\textbf{Korrigierter Code:} Normierung liefert $\vec{a}' = (0{,}5;\ 0{,}5\sqrt{3})$ — exakt derselbe normierte Vektor wie zuvor. Ergebnis: $0{,}5\cdot 0{,}5 + 0{,}5\cdot 0{,}5\sqrt{3} = 0{,}25(1+\sqrt{3}) \approx 0{,}683$ — unverändert.


Schluss: Der Original-Code verletzt die \textbf{Längeninvarianz} der Richtungsableitung: ein Verdoppeln des Eingabevektors verdoppelt fälschlich das Ergebnis. Die mathematische Definition fordert aber, dass die Richtungsableitung nur von der \textbf{Richtung}, nicht von der \textbf{Länge} des Eingabevektors abhängt.



\noindent\textcolor{rulegray}{\rule{\linewidth}{0.4pt}}


\paragraph{Aufgabe 6 — Operator-Identität \& Laplace \textit{(10 Punkte)}}\mbox{}\\[2pt]

a) (3 P) Erkläre die Identität $\operatorname{rot}(\operatorname{grad}(f)) = 0$ und benenne den zugrunde liegenden Satz. Welche unmittelbare Konsequenz hat das für jedes Gradientenfeld?


b) (4 P) Sei $f(x,y) = \dfrac{x^2}{2y}$. Berechne $\nabla f$ und den Laplace-Operator $\Delta f = \operatorname{div}(\operatorname{grad}(f))$.


c) (3 P) Wozu wird der diskretisierte Laplace-Operator in der Bildverarbeitung eingesetzt? Erkläre kurz den Zusammenhang zur zweiten Ableitung.


\textbf{Lösung:}


a) Für $f \in C^2$ gilt $\partial^2 f / \partial x \partial y = \partial^2 f / \partial y \partial x$ (\textbf{Satz von Schwarz}). Damit ist

\[
\operatorname{rot}(\nabla f) = \frac{\partial^2 f}{\partial x \partial y} - \frac{\partial^2 f}{\partial y \partial x} = 0
\]

Konsequenz: \textbf{Jedes Gradientenfeld ist wirbelfrei}. Umgekehrt besitzt jedes wirbelfreie Feld auf einfach zusammenhängenden Gebieten ein Potential.


b) $\nabla f = \begin{pmatrix} \partial f / \partial x \\ \partial f / \partial y \end{pmatrix} = \begin{pmatrix} x/y \\ -x^2/(2y^2) \end{pmatrix}$.


Zweite Ableitungen:

\[
\frac{\partial^2 f}{\partial x^2} = \frac{\partial}{\partial x}\left(\frac{x}{y}\right) = \frac{1}{y}
\]

\[
\frac{\partial^2 f}{\partial y^2} = \frac{\partial}{\partial y}\left(-\frac{x^2}{2y^2}\right) = -\frac{x^2}{2}\cdot(-2)y^{-3} = \frac{x^2}{y^3}
\]

\[
\Delta f = \frac{1}{y} + \frac{x^2}{y^3}
\]


c) Der diskretisierte Laplace-Operator wird in der Bildverarbeitung zur \textbf{Kantenerkennung} eingesetzt. Hintergrund: $\Delta f$ summiert die zweiten Ableitungen — er reagiert daher stark auf Stellen mit abrupten Änderungen der Helligkeit (Nulldurchgänge der zweiten Ableitung markieren Kanten). In flachen Bildregionen ist $\Delta f \approx 0$, an Kanten/Konturen erreicht er hohe Beträge. So kann man Konturen extrahieren, ohne die Kantenrichtung explizit zu kennen.





% ──────────────────────────────────────────────────────────────

\end{document}
